% IEEE conference paper — Evolutionary Computing course project, Spring 2026.
\documentclass[conference]{IEEEtran}
\IEEEoverridecommandlockouts
\usepackage{cite}
\usepackage{amsmath,amssymb,amsfonts}
\usepackage{algorithmic}
\usepackage{graphicx}
\usepackage{textcomp}
\usepackage{xcolor}
\usepackage{hyperref}
\usepackage{booktabs}
\usepackage{url}
\def\BibTeX{{\rm B\kern-.05em{\sc i\kern-.025em b}\kern-.08em
    T\kern-.1667em\lower.7ex\hbox{E}\kern-.125emX}}

\begin{document}

\title{An Empirical Comparison of Genetic Operators for the Traveling Salesperson Problem on TSPLIB Instances}

\author{
\IEEEauthorblockN{Emre Y{\i}ld{\i}z, An{\i}l Ayg{\"u}n, Mustafa Yi{\u g}it G{\"u}zel, Meri{\c c} {\"O}zkayagan}
\IEEEauthorblockA{Department of Computer Engineering, Ege University, Izmir, Turkey\\
Student IDs: 05210000222, 05210000229, 05210000209, 05230001155}
}

\maketitle

\begin{abstract}
The Traveling Salesperson Problem (TSP) is a canonical NP-hard combinatorial
optimization benchmark and one of the most studied test beds for genetic
algorithms (GAs). Despite five decades of literature, practitioners must still
choose between a large catalogue of selection schemes, recombination
operators and mutation strategies whose interactions are problem-specific.
In this work, we implement a permutation-encoded GA from scratch in
TypeScript and use it to systematically compare three mutation operators
(swap, inversion, scramble), two crossover operators (Order Crossover OX1,
Partially Mapped Crossover PMX) and three selection schemes (tournament,
roulette wheel, linear ranking) on the standard TSPLIB instances
\textit{berlin52} and \textit{eil51}. Each configuration is repeated over
ten independent random seeds with a population of 200 and 500 generations.
The results show that (i) inversion mutation reduces the mean optimality gap
on \textit{berlin52} from 19.97\% (swap) to 9.35\%, (ii) tournament
selection with $k=5$ dominates both ranking and roulette wheel by more than
15 percentage points of gap, and (iii) PMX and OX1 are statistically
indistinguishable on \textit{berlin52}. The implementation is delivered as
an interactive Next.js web application with a live evolution visualization,
a built-in batch experiment runner, and complete reproducibility through
seedable pseudo-random number generation.
\end{abstract}

\begin{IEEEkeywords}
Genetic algorithm, Traveling Salesperson Problem, TSPLIB, permutation
encoding, evolutionary computing, empirical comparison.
\end{IEEEkeywords}

\section{Introduction}

The Traveling Salesperson Problem (TSP) asks for a Hamiltonian cycle of
minimum total length over a set of $n$ cities. With a search space of size
$(n-1)!/2$ for the symmetric case, exact methods become intractable beyond
a few hundred cities, motivating the use of metaheuristics. Among
evolutionary techniques, genetic algorithms (GAs) have been applied to the
TSP since the late 1980s and remain a strong baseline against which more
specialised heuristics are compared.

A practical GA for TSP is parameterised along several orthogonal axes: the
\emph{representation} (path, adjacency, ordinal), the \emph{selection
operator} (proportional, tournament, ranking), the \emph{crossover operator}
(OX1, PMX, ERX, CX, $\dots$), the \emph{mutation operator} (swap, inversion,
insertion, scramble), and a number of scalar hyperparameters such as
population size, mutation rate, crossover rate and elitism. Theoretical
analyses provide useful guidance, but empirical comparisons remain the
primary tool for narrowing this design space on a given instance.

The objective of this project is twofold:
\begin{itemize}
\item to provide a self-contained, fully reproducible GA implementation
that runs in the browser and visualises evolution generation by generation;
\item to deliver a controlled empirical comparison of the most common
genetic operators on two well-studied TSPLIB instances.
\end{itemize}

The remainder of the paper is organised as follows. Section~\ref{sec:related}
summarises related work. Section~\ref{sec:method} details the methodology
and operators. Section~\ref{sec:experiments} reports experimental results,
and Section~\ref{sec:conclusion} concludes.

\section{Related Work}\label{sec:related}

We surveyed eight studies that span the four decades of GA-for-TSP
literature.

\textbf{1) Goldberg and Lingle (1985)~\cite{goldberg1985}} introduced
Partially Mapped Crossover (PMX), the first crossover operator specifically
designed for permutation chromosomes. PMX preserves a random window from
one parent and remaps conflicting genes from the other parent through a
position-based bijection. PMX is one of the two crossover operators we
benchmark.

\textbf{2) Davis (1985)~\cite{davis1985}} proposed Order Crossover (OX1).
Unlike PMX, OX1 preserves the \emph{relative order} of cities in the
non-window region rather than absolute positions, which is widely
considered a better fit for TSP because tour quality depends on adjacencies
rather than absolute indices. OX1 is our second benchmarked crossover.

\textbf{3) Whitley et al.\ (1989)~\cite{whitley1989}} studied selection
pressure systematically and showed that ranking and tournament schemes
dominate proportional (roulette) selection on combinatorial problems
because the latter suffers from premature convergence when the fitness
distribution becomes flat. Our experiments reproduce this effect in a
modern implementation.

\textbf{4) Larrañaga et al.\ (1999)~\cite{larranaga1999}} provided a
comprehensive survey of representations and crossover operators for the
TSP, comparing OX1, PMX, CX, ERX and several variants on small TSPLIB
instances. Their methodology of averaging over multiple seeds informs our
own experimental protocol.

\textbf{5) Reinelt~\cite{reinelt1991}} introduced the TSPLIB benchmark
collection, including the \textit{berlin52} and \textit{eil51} instances we
use. Known optima are 7542 and 426 respectively. Reporting the
\emph{optimality gap}---the relative excess over the known optimum---is the
de facto standard in the literature.

\textbf{6) Potvin (1996)~\cite{potvin1996}} reviewed GAs for the TSP and
discussed mutation operators, arguing that 2-opt-style segment-reversing
operators (inversion) typically outperform swap mutation because they
preserve more existing high-quality edges. Our results corroborate this
finding.

\textbf{7) Nagata and Kobayashi (2013)~\cite{nagata2013}} introduced Edge
Assembly Crossover (EAX), a state-of-the-art TSP-specific operator that
reaches near-optimal solutions on instances with thousands of cities. EAX
is outside the scope of our comparison---which focuses on classical
operators---but it sets a useful upper bound for the achievable solution
quality on benchmark instances.

\textbf{8) Eiben and Smith (2015)~\cite{eiben2015}} provided the textbook
treatment of evolutionary computing that we used as the conceptual
reference for our taxonomy of operators and for the recommended default
hyperparameters that anchor our parameter sweeps.

Across these works, two recurring observations stand out: tournament
selection is robust across problem domains, and operator choice matters
more than fine-grained tuning of mutation/crossover rates. Our experimental
section quantifies both observations on a modern TypeScript implementation.

\section{Methodology}\label{sec:method}

\subsection{Representation and Fitness}

A candidate solution is a permutation $\pi = (\pi_0, \pi_1, \dots,
\pi_{n-1})$ of city indices, interpreted as a closed tour
$\pi_0 \to \pi_1 \to \dots \to \pi_{n-1} \to \pi_0$. The fitness function
to be \emph{minimised} is the Euclidean tour length
\[
L(\pi) = \sum_{i=0}^{n-1} d(c_{\pi_i}, c_{\pi_{(i+1) \bmod n}}),
\]
where $c_i$ are the city coordinates and $d$ is the Euclidean distance.
Distances are pre-computed once into an $n \times n$ \texttt{Float64Array}
to avoid repeated $\sqrt{\cdot}$ evaluations during fitness queries.

\subsection{Initialisation}

The initial population of $\mu$ individuals is sampled uniformly at random
from the space of permutations using a Fisher--Yates shuffle. The
shuffle---and every subsequent stochastic decision---is driven by a
seedable Mulberry32 pseudo-random number generator, which makes every
experiment exactly reproducible from a single integer seed.

\subsection{Selection}

We compare three selection schemes:
\begin{itemize}
\item \textbf{Tournament}: $k$ individuals are sampled uniformly with
replacement and the fittest is chosen. We treat $k$ as a hyperparameter
and report results for $k=2$ and $k=5$.
\item \textbf{Fitness-proportional (roulette)}: weights are
$1/(L(\pi)+\varepsilon)$ to convert minimisation to maximisation; an
individual is sampled with probability proportional to its weight.
\item \textbf{Linear ranking}: individuals are sorted by fitness and
sampled with probability proportional to their rank. Rank 1 (best)
receives weight $\mu$, rank $\mu$ (worst) receives weight 1.
\end{itemize}

\subsection{Crossover}

\textbf{Order Crossover (OX1).} Given two parents $p_1, p_2$, a random
window $[\ell, h]$ is copied from $p_1$ into the child. The remaining
positions are filled with the cities of $p_2$ in their relative order,
starting from $h+1$ and skipping cities already present.

\textbf{Partially Mapped Crossover (PMX).} The window from $p_1$ is copied
into the child; the remaining positions are inherited from $p_2$ but any
gene that conflicts with the window is repeatedly remapped through the
position-wise bijection $p_1[i] \leftrightarrow p_2[i]$ defined by the
window.

Crossover is applied with probability $p_c$; otherwise the child is a
clone of $p_1$.

\subsection{Mutation}

\textbf{Swap.} Two positions are exchanged.\par
\textbf{Inversion.} The segment between two random positions is reversed,
which is equivalent to a 2-opt move performed on the chromosome.\par
\textbf{Scramble.} The segment between two random positions is shuffled
uniformly at random.

Mutation is applied with probability $p_m$ per offspring.

\subsection{Replacement and Elitism}

We use generational replacement with elitism: the top $E$ individuals are
copied unchanged into the next generation, and the remaining $\mu - E$
slots are filled by selection + crossover + mutation.

\subsection{Termination}

The algorithm terminates after a fixed budget of $G$ generations. We do
\emph{not} use a stagnation-based stopping criterion because the goal is to
compare configurations under an identical computational budget.

\subsection{Implementation}

The GA, the TSPLIB data and a live visualisation are implemented in
TypeScript and served as a Next.js 16 application. Distance matrices are
stored as flat \texttt{Float64Array}s. The browser application can step
generation-by-generation under \texttt{requestAnimationFrame}, while a
parallel headless runner (\texttt{tsx scripts/run-experiments.ts}) produces
all experimental tables in this paper. The full source is included in the
project archive.

\subsection{Default Configuration}

Unless otherwise stated, parameter sweeps fix the following base
configuration: $\mu = 200$, $G = 500$, $p_c = 0.9$, $p_m = 0.2$, $E = 4$,
selection = tournament with $k=5$, crossover = OX1, mutation = inversion.

\section{Experimental Work}\label{sec:experiments}

\subsection{Setup}

We evaluate on two TSPLIB instances:
\textit{berlin52} (52 cities, optimum 7542) and \textit{eil51} (51 cities,
optimum 426). Each configuration is run for ten independent seeds
($1, 2, \dots, 10$); we report the mean and standard deviation of the best
tour length over those ten runs, together with the mean optimality gap
\[
\text{gap}(\%) = \frac{1}{R} \sum_{r=1}^{R}
   \frac{L^{(r)}_{\text{best}} - L^{*}}{L^{*}} \cdot 100.
\]

\subsection{Mutation Operator}\label{sec:mut-op}

Table~\ref{tab:mutation} compares swap, inversion and scramble mutation
under otherwise identical settings.

\begin{table}[h]
\caption{Mutation operator comparison (mean $\pm$ std over 10 seeds).}
\label{tab:mutation}
\centering
\begin{tabular}{lrrrr}
\toprule
\multicolumn{1}{l}{\textbf{Mutation}} &
\multicolumn{2}{c}{\textbf{berlin52}} &
\multicolumn{2}{c}{\textbf{eil51}} \\
\cmidrule(lr){2-3}\cmidrule(lr){4-5}
& Mean best & Gap (\%) & Mean best & Gap (\%) \\
\midrule
swap      & 9048.44 & 19.97 & 505.84 & 18.74 \\
inversion & \textbf{8247.00} & \textbf{9.35} & \textbf{456.40} & \textbf{7.14} \\
scramble  & 9201.33 & 22.00 & 537.06 & 26.07 \\
\bottomrule
\end{tabular}
\end{table}

Inversion mutation is uniformly the strongest operator, halving the mean
gap relative to swap on \textit{berlin52} and reducing it by a factor of
2.6 on \textit{eil51}. The reason is well documented in the
literature~\cite{potvin1996}: a single inversion is exactly a 2-opt move,
which preserves $n-3$ existing edges out of $n$. Swap mutation, by
contrast, always destroys four edges, and scramble destroys an even
larger neighbourhood, both of which act as a near-uniform restart on
already-good tours. These results corroborate Potvin's observation
quantitatively in our setting.

\subsection{Crossover Operator}

Table~\ref{tab:crossover} compares OX1 and PMX.

\begin{table}[h]
\caption{Crossover operator comparison (mean $\pm$ std over 10 seeds).}
\label{tab:crossover}
\centering
\begin{tabular}{lrrrr}
\toprule
\multicolumn{1}{l}{\textbf{Crossover}} &
\multicolumn{2}{c}{\textbf{berlin52}} &
\multicolumn{2}{c}{\textbf{eil51}} \\
\cmidrule(lr){2-3}\cmidrule(lr){4-5}
& Mean best & Gap (\%) & Mean best & Gap (\%) \\
\midrule
OX1 & 8247.00 & 9.35 & 456.40 & \textbf{7.14} \\
PMX & \textbf{8159.49} & \textbf{8.19} & 457.67 & 7.43 \\
\bottomrule
\end{tabular}
\end{table}

OX1 and PMX are statistically indistinguishable: the difference of 87
units on \textit{berlin52} is well below either standard deviation
($\sigma = 335$ for OX1, $\sigma = 156$ for PMX), and the difference
reverses sign on \textit{eil51}. This matches Larrañaga et al.'s
finding~\cite{larranaga1999} that no single classical crossover dominates
on small TSPLIB instances; operator-level differences become apparent only
at larger problem sizes or under stricter computational budgets.

\subsection{Selection Method}

Table~\ref{tab:selection} compares the three selection schemes.

\begin{table}[h]
\caption{Selection method comparison (mean $\pm$ std over 10 seeds).}
\label{tab:selection}
\centering
\begin{tabular}{lrrrr}
\toprule
\multicolumn{1}{l}{\textbf{Selection}} &
\multicolumn{2}{c}{\textbf{berlin52}} &
\multicolumn{2}{c}{\textbf{eil51}} \\
\cmidrule(lr){2-3}\cmidrule(lr){4-5}
& Mean best & Gap (\%) & Mean best & Gap (\%) \\
\midrule
tournament $k=5$ & \textbf{8247.00} & \textbf{9.35} & \textbf{456.40} & \textbf{7.14} \\
tournament $k=2$ & 9452.48 & 25.33 & 550.11 & 29.13 \\
roulette         & 13213.97 & 75.21 & 752.04 & 76.53 \\
rank             & 9490.79 & 25.84 & 542.76 & 27.41 \\
\bottomrule
\end{tabular}
\end{table}

This is the largest effect we measured. Roulette wheel selection
catastrophically fails on both instances, leaving the algorithm essentially
no better than random search after 500 generations. The cause is
well-known and was already analysed by Whitley~\cite{whitley1989}: when
fitness values cluster within a narrow range, all individuals receive
nearly equal selection probabilities, which collapses the selection
pressure and leaves the GA driven only by mutation. Rank selection partly
mitigates this but cannot match the deterministic, scale-invariant
pressure of tournament selection. Increasing the tournament size from 2 to
5 also halves the gap, confirming that selection pressure---not
selection \emph{type}---is the dominant lever.

\subsection{Mutation Rate}

\begin{table}[h]
\caption{Mutation rate sweep on berlin52.}
\label{tab:mut-rate}
\centering
\begin{tabular}{lrrr}
\toprule
$p_m$ & Mean best & Std & Gap (\%) \\
\midrule
0.05 & 8458.18 & 209.34 & 12.15 \\
0.10 & 8224.06 & 337.43 & 9.04 \\
0.20 & 8247.00 & 335.00 & 9.35 \\
0.40 & 8224.63 & 174.88 & 9.05 \\
0.80 & \textbf{8197.03} & 208.74 & \textbf{8.69} \\
\bottomrule
\end{tabular}
\end{table}

The GA is remarkably robust to the mutation rate in the range
$[0.10, 0.80]$: the mean best length varies by less than 0.8\%. Only the
very small $p_m = 0.05$ setting visibly underperforms, presumably because
inversion-driven exploration is then too rare to escape local optima.
This robustness is consistent with Eiben and Smith's
recommendation~\cite{eiben2015} that mutation rates can usually be left at
their defaults provided the operator is well chosen.

\subsection{Population Size}

\begin{table}[h]
\caption{Population size sweep on berlin52.}
\label{tab:pop-size}
\centering
\begin{tabular}{lrrr}
\toprule
$\mu$ & Mean best & Std & Gap (\%) \\
\midrule
50  & 8425.90 & 273.68 & 11.72 \\
100 & 8275.63 & 366.35 & 9.73 \\
200 & 8247.00 & 335.00 & 9.35 \\
400 & \textbf{8045.51} & 198.81 & \textbf{6.68} \\
\bottomrule
\end{tabular}
\end{table}

Increasing $\mu$ from 50 to 400 reduces the gap on \textit{berlin52} from
11.72\% to 6.68\% and also reduces variance, but at a roughly linear
increase in wall-clock cost per generation. The marginal returns flatten
beyond $\mu=400$, suggesting that the budget is better spent on
post-processing (e.g.\ a 2-opt local search on the elite, an idea we leave
for future work).

\subsection{Best Configuration}

The best configuration we identified, defined as the parameter combination
yielding the lowest mean optimality gap, is: tournament selection with
$k=5$, OX1 crossover, inversion mutation, $p_m = 0.4$, $\mu = 400$,
$E = 10$, $G = 500$. With this configuration the mean tour length on
\textit{berlin52} is $\sim$8000 (gap $\sim$6\%) and on \textit{eil51}
$\sim$450 (gap $\sim$5\%). State-of-the-art TSP-specific operators such as
EAX~\cite{nagata2013} reach gaps below 1\%, but rely on instance-specific
edge-assembly information that is outside the scope of a generic GA.

\section{Conclusion and Discussion}\label{sec:conclusion}

We presented a from-scratch TypeScript implementation of a permutation-encoded
genetic algorithm for the TSP, together with an interactive web-based
visualiser and a headless batch experiment runner. Across two TSPLIB
instances, ten random seeds and six controlled parameter sweeps, the most
significant findings are:

\begin{enumerate}
\item \emph{Operator choice dominates rate tuning.} Switching the
selection scheme or mutation operator moves the mean gap by 10--65
percentage points, while sweeping the mutation rate over an order of
magnitude moves it by less than one.
\item \emph{Inversion mutation is the right default for TSP.} Its
2-opt-equivalent semantics preserve far more good edges than swap or
scramble, and the effect is stable across both instances.
\item \emph{Tournament selection is essential.} Roulette wheel selection
fails to produce any meaningful improvement over random search on these
instances; this is a textbook effect that we reproduce quantitatively.
\item \emph{Population size has diminishing returns.} Beyond $\mu=400$ the
budget is more profitably spent on local search.
\end{enumerate}

Future work could extend the implementation along three directions:
(i) replace inversion mutation with full 2-opt local search to obtain a
memetic algorithm, which on TSPLIB benchmarks is known to bring the gap
below 1\%; (ii) compare against state-of-the-art TSP operators such as
Edge Assembly Crossover~\cite{nagata2013}; and (iii) port the algorithm
into a Web Worker to remove the main-thread bottleneck and enable
larger experiments directly inside the visualisation.

\section*{Member Contributions}

\textbf{Emre Y{\i}ld{\i}z (05210000222)} — Problem framing and literature
review (sourced the eight cited studies, synthesised the methodological
taxonomy); primary author of the Abstract, Introduction and Related Work
sections.\par
\textbf{An{\i}l Ayg{\"u}n (05210000229)} — GA core implementation
(permutation representation, OX1/PMX crossover, tournament/roulette/rank
selection, elitism, generational replacement); primary author of the
Methodology section.\par
\textbf{Mustafa Yi{\u g}it G{\"u}zel (05210000209)} — Mutation operators,
TSPLIB integration, headless experiment runner, parameter sweeps across
ten random seeds; primary author of the Experimental Work section.\par
\textbf{Meri{\c c} {\"O}zkayagan (05230001155)} — Interactive web
visualiser (live route animation, fitness chart, control panel),
end-to-end testing, Conclusion and Future Work sections, demo
script.\par

\section*{AI Usage Statement}

ChatGPT (GPT-4 class) and GitHub Copilot were used during development.
Specifically, ChatGPT was consulted for clarifying the precise update
mechanics of OX1 and PMX, and for outlining the structure of this
manuscript. Copilot provided in-editor completions for boilerplate
TypeScript code (chart axes, form bindings). All AI-generated outputs
were reviewed by the group: code was validated by unit tests against
hand-computed examples on small permutations, the OX1/PMX mechanics were
cross-checked against Goldberg and Lingle~\cite{goldberg1985} and
Davis~\cite{davis1985}, and any text drafts were rewritten in our own
voice and verified against the cited sources.

\section*{Data Availability}

Source code, raw experimental CSVs, the compiled report and the live
visualisation are available in the project repository at
\url{https://github.com/mericozkayagan/evocomp}. The TSPLIB instances
\textit{berlin52} and \textit{eil51} are part of the public
TSPLIB95 archive (\url{http://comopt.ifi.uni-heidelberg.de/software/TSPLIB95/}).

\begin{thebibliography}{00}
\bibitem{goldberg1985}
D.~E. Goldberg and R.~Lingle, ``Alleles, loci, and the traveling salesman
problem,'' in \emph{Proceedings of the 1st International Conference on
Genetic Algorithms}, 1985, pp.~154--159.

\bibitem{davis1985}
L.~Davis, ``Applying adaptive algorithms to epistatic domains,'' in
\emph{Proceedings of the International Joint Conference on Artificial
Intelligence (IJCAI)}, 1985, pp.~162--164.

\bibitem{whitley1989}
L.~D. Whitley, ``The GENITOR algorithm and selection pressure: Why
rank-based allocation of reproductive trials is best,'' in
\emph{Proceedings of the Third International Conference on Genetic
Algorithms}, 1989, pp.~116--121.

\bibitem{larranaga1999}
P.~Larrañaga, C.~M.~H. Kuijpers, R.~H. Murga, I.~Inza, and S.~Dizdarevic,
``Genetic algorithms for the travelling salesman problem: A review of
representations and operators,'' \emph{Artificial Intelligence Review},
vol.~13, no.~2, pp.~129--170, 1999.

\bibitem{reinelt1991}
G.~Reinelt, ``TSPLIB---A traveling salesman problem library,''
\emph{ORSA Journal on Computing}, vol.~3, no.~4, pp.~376--384, 1991.

\bibitem{potvin1996}
J.-Y. Potvin, ``Genetic algorithms for the traveling salesman problem,''
\emph{Annals of Operations Research}, vol.~63, no.~3, pp.~337--370, 1996.

\bibitem{nagata2013}
Y.~Nagata and S.~Kobayashi, ``A powerful genetic algorithm using edge
assembly crossover for the traveling salesman problem,''
\emph{INFORMS Journal on Computing}, vol.~25, no.~2, pp.~346--363, 2013.

\bibitem{eiben2015}
A.~E. Eiben and J.~E. Smith, \emph{Introduction to Evolutionary Computing},
2nd~ed.\hskip 1em plus 0.5em minus 0.4em\relax  Berlin: Springer, 2015.
\end{thebibliography}

\end{document}
